\makeatletter
\@ifundefined{keepfig}{\newenvironment{keepfig}{\par\edef\keepfig@pi{\the\parindent}\noindent
  \begin{minipage}{\linewidth}\parindent=\keepfig@pi\relax}{\end{minipage}\par}}{}
\makeatother
\section{Основная теорема теории Галуа}

Пусть $E$~--- расширение поля $F$. \emph{Подрасширением} расширения $E/F$ называется расширение $M/F$ с $M\subset E$,
т.\,е.\ поле $M$, для которого $F\subset M\subset E$ (говорят также: \emph{промежуточное поле}). Если $E/F$~--- расширение
Галуа, то подрасширения взаимно однозначно соответствуют подгруппам группы $\Gal(E/F)$. Это соответствие позволяет
переводить вопросы о полях (обычно трудные) в вопросы о конечных группах (обычно гораздо более лёгкие); например,
у расширения Галуа лишь конечное число подрасширений, и все они находятся перебором подгрупп.

Все ингредиенты уже готовы: теорема Артина (\rf{3.4}, \rf{3.5}, \rf{3.11}) говорит, что для конечной группы $H$
автоморфизмов поля $E$ расширение $E/E^H$ является расширением Галуа с группой ровно $H$, а \rf{3.10} и \rf{3.13}~--- что
для расширения Галуа $E/F$ и промежуточного поля $M$ расширение $E/M$ тоже Галуа, т.\,е.\ $M$ есть неподвижное поле
группы $\Gal(E/M)$. Остаётся аккуратно сопоставить эти факты.

\begin{theorem}[3.17: основная теорема теории Галуа]
Пусть $E/F$~--- расширение Галуа с группой $G=\Gal(E/F)$. Отображение $H\mapsto E^H$~--- биекция множества подгрупп $G$
на множество подрасширений $E/F$,
\[
\{\text{подгруппы } H\subset G\}\ \overset{1:1}{\longleftrightarrow}\ \{\text{подрасширения }F\subset M\subset E\},
\]
и обратное к нему отображение~--- $M\mapsto\Gal(E/M)$. Кроме того:
\begin{enumerate}[label=\textup{(\alph*)}]
\item $H_1\supset H_2\iff E^{H_1}\subset E^{H_2}$ (соответствие обращает включения);
\item $(H_1:H_2)=[E^{H_2}:E^{H_1}]$ при $H_1\supset H_2$;
\item $\sigma H\sigma^{-1}\leftrightarrow\sigma M$ для $\sigma\in G$, т.\,е.
\[
E^{\sigma H\sigma^{-1}}=\sigma(E^H),\qquad \Gal(E/\sigma M)=\sigma\Gal(E/M)\sigma^{-1};
\]
\item подгруппа $H$ нормальна в $G$ тогда и только тогда, когда $E^H$ нормально (а значит, является расширением Галуа)
над $F$; в этом случае
\[
\Gal(E^H/F)\cong G/H .
\]
\end{enumerate}
\end{theorem}

Здесь $(H_1:H_2)$~--- индекс подгруппы $H_2$ в $H_1$, т.\,е.\ число смежных классов $hH_2$, $h\in H_1$; по теореме Лагранжа
$(H_1:H_2)=|H_1|/|H_2|$. В частности, $[E:E^H]=|H|$ и $[E^H:F]=(G:H)$. Например, $G=1$ лишь при $E=F$.

\begin{proof}
\emph{Биекция.} Каждая подгруппа $H\subset G$~--- конечная группа автоморфизмов поля $E$, поэтому по теореме Артина
(\rf{3.11}) расширение $E/E^H$~--- расширение Галуа, причём $\Gal(E/E^H)=H$ и $[E:E^H]=|H|$. Кроме того, $F=E^G\subset
E^H$, так что $E^H$~--- подрасширение. Значит, композиция $H\mapsto E^H\mapsto\Gal(E/E^H)$ тождественна. Обратно, пусть
$M$~--- подрасширение. Каждый автоморфизм $E$, тождественный на $M$, тождествен и на $F\subset M$, поэтому
$\Gal(E/M)$~--- подгруппа $G$. По \rf{3.13} расширение $E/M$ является расширением Галуа, т.\,е.\ (\rf{3.10}, (d)$\Rightarrow$(b))
$E^{\Gal(E/M)}=M$: композиция $M\mapsto\Gal(E/M)\mapsto E^{\Gal(E/M)}$ тоже тождественна. Итак, два отображения взаимно
обратны.

(a) Очевидны импликации
\[
H_1\supset H_2\ \Longrightarrow\ E^{H_1}\subset E^{H_2}\qquad\text{и}\qquad M_1\subset M_2\ \Longrightarrow\ \Gal(E/M_1)\supset\Gal(E/M_2)
\]
(элемент, неподвижный относительно бóльшей группы, неподвижен и относительно меньшей; автоморфизм, тождественный на
бóльшем поле, тождествен и на меньшем). Применяя вторую импликацию к $M_i=E^{H_i}$ и пользуясь равенствами
$\Gal(E/E^{H_i})=H_i$, получаем (a).

(b) По доказанному $|H|=[E:E^H]$ для каждой подгруппы $H$. При $H_1\supset H_2$ имеем $|H_1|=(H_1:H_2)\,|H_2|$ (Лагранж) и
$[E:E^{H_1}]=[E:E^{H_2}]\,[E^{H_2}:E^{H_1}]$ (\rf{1.20}); сравнивая, получаем $(H_1:H_2)=[E^{H_2}:E^{H_1}]$.

(c) Пусть $\sigma\in G$ и $M$~--- подрасширение. Поле $\sigma M$ содержит $\sigma F=F$, т.\,е.\ тоже подрасширение. Для
$\tau\in G$
\begin{multline*}
\tau\in\Gal(E/\sigma M)\iff\tau(\sigma\alpha)=\sigma\alpha\ \ \forall\alpha\in M\\
\iff(\sigma^{-1}\tau\sigma)(\alpha)=\alpha\ \ \forall\alpha\in M\iff\sigma^{-1}\tau\sigma\in\Gal(E/M),
\end{multline*}
а последнее равносильно $\tau\in\sigma\Gal(E/M)\sigma^{-1}$. Это второе равенство. Первое проверяется так же: для
$\alpha\in E$
\begin{multline*}
\alpha\in E^{\sigma H\sigma^{-1}}\iff\sigma\tau\sigma^{-1}(\alpha)=\alpha\ \ \forall\tau\in H\\
\iff\tau(\sigma^{-1}\alpha)=\sigma^{-1}\alpha\ \ \forall\tau\in H\iff\sigma^{-1}\alpha\in E^H\iff\alpha\in\sigma(E^H).
\end{multline*}

(d) Пусть $M=E^H$, так что $H=\Gal(E/M)$. Покажем, что следующие условия равносильны:
\[
\text{(i) } H\triangleleft G;\qquad\text{(ii) } \sigma M=M\text{ для всех }\sigma\in G;\qquad\text{(iii) } M/F\text{ нормально.}
\]
В (ii) имеется в виду равенство множеств: $\sigma$ переводит $M$ в себя, но, вообще говоря, двигает его элементы
(поэлементно $M$ неподвижно лишь относительно $H$); так же в (i) $\sigma H\sigma^{-1}=H$~--- равенство подгрупп, а не
$\sigma\tau\sigma^{-1}=\tau$ для каждого $\tau\in H$.

(i)$\Leftrightarrow$(ii). По (c) и биективности соответствия $\sigma H\sigma^{-1}=H\iff E^{\sigma H\sigma^{-1}}=E^H\iff\sigma M=M$.

(iii)$\Rightarrow$(ii). Расширение $M/F$ конечно, $M=F[\alpha_1,\dots,\alpha_m]$. Пусть $\mu_i\in F[X]$~--- минимальный
многочлен $\alpha_i$ над $F$ и $\sigma\in G$. Применяя $\sigma$ к равенству $\mu_i(\alpha_i)=0$ (коэффициенты $\mu_i$ лежат
в $F$ и неподвижны), получаем $\mu_i(\sigma\alpha_i)=0$. Так как $M/F$ нормально, $\mu_i=\prod_j(X-\beta_{ij})$ с
$\beta_{ij}\in M$; значит, $\prod_j(\sigma\alpha_i-\beta_{ij})=0$ в поле $E$, и $\sigma\alpha_i$ равен одному из
$\beta_{ij}\in M$. Поэтому $\sigma M=F[\sigma\alpha_1,\dots,\sigma\alpha_m]\subset M$; а так как $\sigma$ инъективно и
$F$-линейно, $[\sigma M:F]=[M:F]$, и $\sigma M=M$.

(ii)$\Rightarrow$(iii) и изоморфизм $\Gal(M/F)\cong G/H$. Если $\sigma M=M$ для всех $\sigma\in G$, то ограничение
$\sigma|_M$~--- автоморфизм поля $M$ над $F$, и мы получаем гомоморфизм групп
\[
r\colon G\to\Aut(M/F),\qquad \sigma\mapsto\sigma|_M .
\]
Его ядро состоит из $\sigma$, тождественных на $M$, т.\,е.\ равно $\Gal(E/M)=H$. Поэтому $r$ индуцирует инъективный
гомоморфизм $G/H\to\Aut(M/F)$, и его образ $\bar G=r(G)$~--- конечная группа автоморфизмов поля $M$, изоморфная $G/H$.
Найдём её неподвижное поле: для $\alpha\in M$
\[
\alpha\in M^{\bar G}\iff\sigma\alpha=\alpha\ \ \forall\sigma\in G\iff\alpha\in E^G=F
\]
(последнее равенство~--- \rf{3.10}(b), так как $E/F$~--- расширение Галуа). Итак, $M^{\bar G}=F$. По теореме Артина
(\rf{3.11}), применённой к полю $M$ и группе $\bar G$, расширение $M/F$~--- расширение Галуа (в частности, нормально), и
$\Gal(M/F)=\bar G\cong G/H$.

Наконец, нормальное подрасширение $M$ автоматически является расширением Галуа: оно конечно, а сепарабельно потому,
что минимальный многочлен над $F$ каждого элемента $M\subset E$ сепарабелен ($E/F$ сепарабельно).
\end{proof}

Подчеркнём, в каком порядке шло рассуждение в (d): мы \emph{не} пользовались тем, что $M/F$~--- расширение Галуа (это как
раз доказывается). Мы построили конкретную конечную группу $\bar G$ автоморфизмов поля $M$, проверили, что её неподвижное
поле равно $F$ (это сводится к $E^G=F$), и лишь затем применили теорему Артина. Равенство $(E^H)^{G/H}=F$, в котором
$G/H$ действует на $E^H$ через ограничения, означает именно $M^{\bar G}=F$.

\begin{keepfig}
\begin{center}
\begin{tikzpicture}[every node/.style={font=\small}]
\begin{scope}
\node[cold] (one) at (0,3.3) {$1$};
\node (H2) at (0,2.2) {$H_2$};
\node (H1) at (0,1.1) {$H_1$};
\node[cold] (G) at (0,0) {$G$};
\draw (one) -- node[left,font=\scriptsize]{$|H_2|$} (H2) -- node[left,font=\scriptsize]{$(H_1:H_2)$} (H1) -- node[left,font=\scriptsize]{$(G:H_1)$} (G);
\node[font=\scriptsize,cold] at (0,-0.6) {подгруппы};
\end{scope}
\begin{scope}[xshift=4.2cm]
\node[cold] (E) at (0,3.3) {$E$};
\node[cnew] (M2) at (0,2.2) {$E^{H_2}$};
\node[cnew] (M1) at (0,1.1) {$E^{H_1}$};
\node[cold] (F) at (0,0) {$F$};
\draw (E) -- node[right,font=\scriptsize]{$|H_2|$} (M2) -- node[right,font=\scriptsize]{$(H_1:H_2)$} (M1) -- node[right,font=\scriptsize]{$(G:H_1)$} (F);
\node[font=\scriptsize,cold] at (0,-0.6) {подрасширения};
\end{scope}
\draw[cnew,->] (0.55,2.2) -- node[above,font=\scriptsize]{$H\mapsto E^H$} (3.5,2.2);
\draw[cnew,<-] (0.55,1.1) -- node[below,font=\scriptsize]{$M\mapsto\Gal(E/M)$} (3.5,1.1);
\end{tikzpicture}
\end{center}
\figcap{Основная теорема: подгруппы рисуем <<вверх ногами>> (вверху~--- меньшие), подполя~--- обычным образом (вверху~---
бóльшие); тогда соответствующие объекты стоят на одной высоте, а индексы на левых рёбрах равны степеням на правых.}
\end{keepfig}

\begin{remark}[3.17a]
Пункт (c) означает, что соответствие Галуа \emph{согласовано с действием} группы $G$: группа $G$ действует на множестве
своих подгрупп сопряжениями ($H\mapsto\sigma H\sigma^{-1}$), а на множестве подрасширений~--- просто как группа
автоморфизмов поля ($M\mapsto\sigma M$), и биекция $H\mapsto E^H$ переводит одно действие в другое. Подгруппы,
неподвижные относительно сопряжений,~--- это нормальные подгруппы; подрасширения, неподвижные относительно $G$,~--- это (по
доказательству (d)) нормальные подрасширения. Так что (d) отчасти является частным случаем (c). Поля $\sigma M$ называются
\emph{сопряжёнными} с $M$; их число равно числу подгрупп, сопряжённых с $H=\Gal(E/M)$, т.\,е.\ индексу нормализатора
$(G:N_G(H))$ (о самой группе $N_G(H)/H$ см.\ упражнение 3-4).
\end{remark}

Сюръективность гомоморфизма $G\to\Gal(E^H/F)$ в (d) можно получить и без теоремы Артина~--- продолжением автоморфизмов.
Это рассуждение даёт больше: оно описывает все вложения произвольного (не обязательно нормального) подрасширения в $E$.

\begin{proposition}[3.17b]
Пусть $E/F$~--- расширение Галуа с группой $G$, $M$~--- подрасширение и $H=\Gal(E/M)$. Каждый $F$-гомоморфизм
$\varphi\colon M\to E$ продолжается до некоторого $\sigma\in G$, и ограничение $\sigma\mapsto\sigma|_M$ задаёт биекцию
\[
G/H\ \xrightarrow{\ \sim\ }\ \Hom_F(M,E),\qquad \sigma H\mapsto\sigma|_M .
\]
В частности, $F$-гомоморфизмов $M\to E$ ровно $(G:H)=[M:F]$, и если $\sigma M=M$ для всех $\sigma\in G$, то
$\Aut(M/F)=\{\sigma|_M\}\cong G/H$.
\end{proposition}

\begin{proof}
По \rf{3.10} $E$~--- поле разложения над $F$ некоторого сепарабельного многочлена $f\in F[X]$; тогда $E$ порождено над $M$
корнями $f$. Пусть $\varphi\colon M\to E$~--- $F$-гомоморфизм. Поскольку коэффициенты $f$ лежат в $F$, имеем $\varphi f=f$, и
$f$ раскладывается в $E$ на линейные множители. По \rf{2.12a} (с $\varphi_0=\varphi$, $\Om=E$) гомоморфизм $\varphi$ продолжается
до гомоморфизма $\sigma\colon E\to E$; он инъективен и $F$-линеен, а $[E:F]<\infty$, поэтому $\sigma$~--- автоморфизм,
$\sigma\in G$. Далее, $\sigma|_M=\sigma'|_M$ $\iff$ $\sigma^{-1}\sigma'$ тождествен на $M$ $\iff$ $\sigma^{-1}\sigma'\in H$
$\iff$ $\sigma H=\sigma'H$. Значит, отображение $\sigma H\mapsto\sigma|_M$ корректно, инъективно и (по доказанному)
сюръективно. Число смежных классов равно $(G:H)=[M:F]$ по \rf{3.17}(b). Последнее утверждение: если $\sigma M=M$ для всех
$\sigma$, то каждый $F$-автоморфизм $M$~--- это $F$-гомоморфизм $M\to E$ с образом $M$, т.\,е.\ имеет вид $\sigma|_M$.
\end{proof}

Отсюда второе доказательство пункта (d), в котором теорема Артина к полю $M=E^H$ не применяется. Равносильность
(i)$\Leftrightarrow$(ii) и импликация (iii)$\Rightarrow$(ii) доказаны выше, и там она к $M$ тоже не
применялась. Пусть выполнено (ii). Для $\alpha\in M$ корни
его минимального многочлена над $F$~--- это сопряжённые $\sigma\alpha$, $\sigma\in G$ (\rf{3.14}), и все они лежат в $\sigma M=M$;
значит, $M/F$ нормально, а будучи конечным и сепарабельным (как подрасширение $E/F$), является расширением Галуа по
определению \rf{3.9}. По \rf{3.17b} гомоморфизм ограничения $G\to\Aut(M/F)$ сюръективен, его ядро~--- $\Gal(E/M)=H$, и
$\Gal(M/F)\cong G/H$. Отметим, что в
\rf{3.17b} не предполагается нормальность $M$: например, для поля разложения $E$ многочлена $X^3-2$ над $\Q$ и
$M=\Q[\sqrt[3]2]$ имеем $(G:H)=3$, и три $\Q$-вложения $M\to E$ переводят $\sqrt[3]2$ в три корня $X^3-2$.

\begin{remark}[3.17c]
Предположение, что $E/F$~--- расширение Галуа, существенно. Пусть $E/F$~--- конечное расширение и $G=\Aut(E/F)$.
Отображение $H\mapsto E^H$ (подгруппы $G$ $\to$ подрасширения) по-прежнему инъективно, так как $\Aut(E/E^H)=H$ (\rf{3.5}), но
его образ состоит лишь из тех $M$, для которых $E/M$~--- расширение Галуа (по \rf{3.11} и \rf{3.10}). Например, у
$\Q[\sqrt[3]2]/\Q$ группа автоморфизмов тривиальна (два других корня $X^3-2$ не вещественны), а подрасширений два: $\Q$ и
$\Q[\sqrt[3]2]$. Так же у чисто несепарабельного расширения $\F_p(T)/\F_p(T^p)$ степени $p$ группа автоморфизмов тривиальна
(\rf{2.3}: корень $X^p-T^p$ единствен), а подрасширений два. Для бесконечных расширений Галуа соответствие сохраняется, если
рассматривать лишь \emph{замкнутые} подгруппы группы Галуа с топологией Крулля (\rf{7.13}). А для трансцендентного
расширения с бесконечной группой автоморфизмов, например $F(X)/F$, аналога основной теоремы нет (\rf{3.18a},
упражнение 3-1).
\end{remark}

\subsection*{Композиты и пересечения}

Напомним, что композит $M_1M_2\cdots M_r$ подполей $M_i$ поля $E$~--- это наименьшее подполе $E$, содержащее все $M_i$; если
$M_i\supset F$ и $M_i/F$ алгебраичны, то $M_1\cdots M_r=F[M_1,\dots,M_r]$: кольцо, порождённое над $F$ алгебраическими
элементами, уже является полем (\rf{1.32}).

\begin{remark}[3.18]
Пусть $E/F$~--- расширение Галуа с группой $G$. Из основной теоремы немедленно получаем следующее.

(a) Пусть $M_1,\dots,M_r$~--- подрасширения и $H_i=\Gal(E/M_i)$. Композит $M_1\cdots M_r$~--- наименьшее поле, содержащее
все $M_i$, и поскольку соответствие обращает включения, ему отвечает наибольшая подгруппа, содержащаяся во всех $H_i$, т.\,е.\
$\bigcap H_i$:
\[
\Gal(E/M_1\cdots M_r)=H_1\cap\dots\cap H_r .
\]
(Можно и напрямую: автоморфизм тождествен на поле, порождённом объединением $\bigcup M_i$, тогда и только тогда, когда он
тождествен на каждом $M_i$.) Двойственно, пересечению полей отвечает подгруппа, \emph{порождённая} подгруппами:
\[
E^{H_1}\cap\dots\cap E^{H_r}=E^{\langle H_1,\dots,H_r\rangle},
\]
так как элемент неподвижен относительно всех элементов каждой $H_i$ тогда и только тогда, когда он неподвижен
относительно всех произведений таких элементов.

(b) Пусть $H$~--- подгруппа $G$ и $M=E^H$. Подгруппа
\[
N=\bigcap_{\sigma\in G}\sigma H\sigma^{-1}
\]
~--- наибольшая нормальная подгруппа $G$, содержащаяся в $H$. Действительно, $N\subset H$ (возьмём $\sigma=1$); для
$\tau\in G$ имеем $\tau N\tau^{-1}=\bigcap_\sigma(\tau\sigma)H(\tau\sigma)^{-1}=N$ (при $\sigma$, пробегающем $G$, произведение
$\tau\sigma$ тоже пробегает $G$); а если $K\triangleleft G$ и $K\subset H$, то $K=\sigma K\sigma^{-1}\subset\sigma H\sigma^{-1}$ для
всех $\sigma$, т.\,е.\ $K\subset N$. По \rf{3.17}(a),(d) нормальные над $F$ подрасширения, содержащие $M$, отвечают
нормальным подгруппам, содержащимся в $H$; значит, $E^N$~--- \emph{наименьшее нормальное над $F$ подрасширение $E$,
содержащее $M$}. По (a) и \rf{3.17}(c)
\[
E^N=E^{\bigcap_\sigma\sigma H\sigma^{-1}}=\prod_{\sigma\in G}E^{\sigma H\sigma^{-1}}=\prod_{\sigma\in G}\sigma M
\]
~--- композит всех полей, сопряжённых с $M$. Поле $E^N$ называется \emph{нормальным замыканием} (или \emph{замыканием
Галуа}) $M$ в $E$; это расширение Галуа поля $F$ (ср.\ \rf{3.12a}(a)).
\end{remark}

Например, если $E$~--- поле разложения $X^3-2$ над $\Q$ и $M=\Q[\sqrt[3]2]$, то $H$~--- подгруппа порядка $2$ в $G\cong S_3$,
пересечение трёх сопряжённых с ней подгрупп тривиально, и нормальное замыкание $M$~--- всё поле $E=\Q[\sqrt[3]2]\cdot\Q[\omega\sqrt[3]2]$.

\begin{aside}[3.18a]
В конечном расширении Галуа степень $E$ над пересечением двух подполей не больше $|G|$: $[E:M_1\cap M_2]=|\langle H_1,H_2\rangle|$.
Для бесконечных групп автоморфизмов это не так: в поле $F(X)$ (характеристика $0$) подполя $F(X^2)$ и $F(X^2-X)$ имеют
степень $2$ (т.\,е.\ $[F(X):F(X^2)]=[F(X):F(X^2-X)]=2$), а их пересечение равно $F$, и $[F(X):F]=\infty$; доказательство (через бесконечную группу,
порождённую двумя автоморфизмами порядка $2$)~--- в решении упражнения 3-1.
\end{aside}

\begin{proposition}[3.19]
Пусть $E$ и $L$~--- расширения $F$, содержащиеся в общем поле $\Om$. Если $E/F$~--- расширение Галуа, то $EL/L$ и
$E/E\cap L$~--- расширения Галуа, и ограничение
\[
\sigma\mapsto\sigma|_E\colon\ \Gal(EL/L)\to\Gal(E/E\cap L)
\]
~--- изоморфизм.
\end{proposition}

Здесь $L$~--- \emph{любое} расширение $F$ внутри $\Om$: оно не обязано быть ни конечным, ни алгебраическим над $F$
(например, $L=F(t)$ или $L=\R$).

\begin{keepfig}
\begin{center}
\begin{tikzpicture}[every node/.style={font=\small}]
\node[cnew] (EL) at (0,2.6) {$EL$};
\node[cold] (E) at (-1.9,1.2) {$E$};
\node[cold] (L) at (1.9,1.2) {$L$};
\node[cnew] (K) at (0,0) {$E\cap L$};
\node[cold] (F) at (0,-1.1) {$F$};
\draw (EL) -- (E);
\draw (EL) -- node[above right,font=\scriptsize,cnew]{$\Gal(EL/L)$} (L);
\draw (E) -- node[above right=-1pt,font=\scriptsize,cnew]{$\Gal(E/E\cap L)$} (K);
\draw (L) -- (K) (K) -- (F);
\draw[cold,dashed] (E) to[bend right=45] node[left,font=\scriptsize]{Галуа} (F);
\node[font=\small] at (5.9,1.2) {$\Gal(EL/L)\ \xrightarrow[\ \sim\ ]{\ \sigma\mapsto\sigma|_E\ }\ \Gal(E/E\cap L)$};
\end{tikzpicture}
\end{center}
\figcap{Предложение \rf{3.19}: при подъёме расширения Галуа $E/F$ на $L$ группа Галуа не меняется, если заменить $F$ на
$E\cap L$; <<параллельные>> стороны ромба $EL/L$ и $E/E\cap L$ имеют одинаковые группы и одинаковые степени.}
\end{keepfig}

\begin{proof}
По \rf{3.10} $E$~--- поле разложения над $F$ сепарабельного многочлена $f\in F[X]$; пусть $\alpha_1,\dots,\alpha_n\in E$~---
его корни, $E=F[\alpha_1,\dots,\alpha_n]$. Тогда $EL=L[\alpha_1,\dots,\alpha_n]$ (это поле содержит $L$ и $E$ и содержится в любом
поле, содержащем $L$ и $E$), т.\,е.\ $EL$~--- поле разложения $f$ над $L$; а $E=(E\cap L)[\alpha_1,\dots,\alpha_n]$~--- поле
разложения $f$ над $E\cap L$. Многочлен $f$ сепарабелен, поэтому по \rf{3.10} расширения $EL/L$ и $E/E\cap L$~---
расширения Галуа.

Каждый $\sigma\in\Gal(EL/L)$ фиксирует коэффициенты $f$ и потому переставляет его корни $\alpha_i$; значит, $\sigma E=E$, и
ограничение $\sigma|_E$~--- автоморфизм $E$, тождественный на $E\cap L\subset L$. Получаем гомоморфизм
$r\colon\Gal(EL/L)\to\Gal(E/E\cap L)$, $r(\sigma)=\sigma|_E$.

\emph{Инъективность.} Если $\sigma|_E=\id$, то $\sigma$ тождествен на $E$ и на $L$, а значит, на поле $EL$, порождённом ими.

\emph{Сюръективность.} Образ $r$~--- подгруппа $\bar G\subset\Gal(E/E\cap L)$. Найдём её неподвижное поле: для $\alpha\in E$
\[
\alpha\in E^{\bar G}\iff\sigma\alpha=\alpha\ \ \forall\sigma\in\Gal(EL/L)\iff\alpha\in(EL)^{\Gal(EL/L)}=L
\]
(последнее равенство~--- \rf{3.10}(b) для расширения Галуа $EL/L$). Значит, $E^{\bar G}=E\cap L=E^{\Gal(E/E\cap L)}$, и по
основной теореме (\rf{3.17}) для расширения Галуа $E/E\cap L$ подгруппы $\bar G$ и $\Gal(E/E\cap L)$ совпадают (иначе:
$\bar G=\Aut(E/E^{\bar G})=\Gal(E/E\cap L)$ по \rf{3.5}).
\end{proof}

\begin{example}[3.19a]
(1) $E=\Q[i]$, $L=\R$ (внутри $\CC$): $E\cap L=\Q$, $EL=\CC$, и \rf{3.19} утверждает, что $\Gal(\CC/\R)\cong\Gal(\Q[i]/\Q)$:
комплексное сопряжение ограничивается до $i\mapsto-i$.

(2) $E=\Q[\zeta]$, $\zeta=e^{2\pi i/7}$, $L=\R$: здесь $E\cap L$~--- неподвижное поле ограничения комплексного сопряжения,
т.\,е.\ поле $\Q[\zeta+\bar\zeta]$ степени $3$ (\rf{3.22}), и $\Gal(\CC/\R)\cong\Gal(\Q[\zeta]/\Q[\zeta+\bar\zeta])$. Вообще,
\rf{3.19} с $L=\R$ показывает, что $E\cap\R$ для расширения Галуа $E\subset\CC$ поля $\Q$~--- неподвижное поле
комплексного сопряжения, и $[E:E\cap\R]\le2$.

(3) Условие <<$E/F$~--- Галуа>> существенно. Пусть $\omega=e^{2\pi i/3}$, $E=\Q[\sqrt[3]2]$, $L=\Q[\omega\sqrt[3]2]$. Это
различные поля степени $3$ (первое вещественно, второе нет), поэтому $E\cap L=\Q$ (степень $[E\cap L:\Q]$ делит $3$ и не
равна $3$). Композит $EL=\Q[\sqrt[3]2,\omega]$~--- поле разложения $X^3-2$, степени $6$. Итак, $[EL:L]=2$, а
$[E:E\cap L]=3$: группы (и степени) <<параллельных сторон>> различны.
\end{example}

\begin{corollary}[3.20]
Пусть в \rf{3.19} расширение $L/F$ конечно. Тогда
\[
[EL:F]=\frac{[E:F]\,[L:F]}{[E\cap L:F]} .
\]
\end{corollary}

\begin{proof}
По \rf{1.20} $[EL:F]=[EL:L]\,[L:F]$. По \rf{3.19} $[EL:L]=|\Gal(EL/L)|=|\Gal(E/E\cap L)|=[E:E\cap L]$, а по \rf{1.20}
$[E:E\cap L]=[E:F]/[E\cap L:F]$.
\end{proof}

\begin{remark}[3.20a]
Без предположения о том, что $E/F$~--- расширение Галуа, верно лишь неравенство: если $E/F$ конечно и $K=E\cap L$, то
\[
[EL:L]\le[E:K],\qquad\text{а при конечном }L/F\quad [EL:F]\le\frac{[E:F]\,[L:F]}{[E\cap L:F]} .
\]
Действительно, пусть $e_1,\dots,e_n$~--- базис $E$ над $K$. Кольцо $L[E]$, порождённое $L$ и $E$, состоит из сумм
произведений $l\,x$ ($l\in L$, $x\in E$), а каждое $x\in E$~--- линейная комбинация $e_j$ с коэффициентами из $K\subset L$;
значит, $e_1,\dots,e_n$ порождают $L[E]$ как векторное пространство над $L$. Конечномерное над полем $L$ кольцо без делителей
нуля является полем (\rf{1.23}), поэтому $L[E]=EL$ и $[EL:L]\le n$. В примере \rf{3.19a}(3) неравенство строгое:
$[EL:\Q]=6<9$.
\end{remark}

\begin{proposition}[3.21]
Пусть $E_1$ и $E_2$~--- расширения $F$, содержащиеся в общем поле. Если $E_1/F$ и $E_2/F$~--- расширения Галуа, то $E_1E_2$
и $E_1\cap E_2$~--- расширения Галуа поля $F$, и отображение
\[
\sigma\mapsto(\sigma|_{E_1},\sigma|_{E_2})\colon\ \Gal(E_1E_2/F)\to\Gal(E_1/F)\times\Gal(E_2/F)
\]
~--- изоморфизм группы $\Gal(E_1E_2/F)$ на подгруппу
\[
H=\{(\sigma_1,\sigma_2)\mid\sigma_1|_{E_1\cap E_2}=\sigma_2|_{E_1\cap E_2}\}.
\]
\end{proposition}

Для гомоморфизмов групп $\pi_1\colon G_1\to G_0$, $\pi_2\colon G_2\to G_0$ подгруппа
$G_1\times_{G_0}G_2=\{(a,b)\in G_1\times G_2\mid\pi_1(a)=\pi_2(b)\}$ называется \emph{расслоённым произведением} $G_1$ и $G_2$
над $G_0$. В этих терминах \rf{3.21} утверждает, что
\[
\Gal(E_1E_2/F)\cong\Gal(E_1/F)\times_{\Gal(E_1\cap E_2/F)}\Gal(E_2/F),
\]
где $\pi_i$~--- ограничения на $E_1\cap E_2$ (они определены, так как $E_1\cap E_2$ нормально над $F$; см.\ доказательство).

\begin{keepfig}
\begin{center}
\begin{tikzpicture}[every node/.style={font=\small}]
\begin{scope}
\node[cnew] (EE) at (0,2.4) {$E_1E_2$};
\node[cold] (E1) at (-1.2,1.2) {$E_1$};
\node[cold] (E2) at (1.2,1.2) {$E_2$};
\node[cnew] (K) at (0,0) {$E_1\cap E_2$};
\node[cold] (F) at (0,-1) {$F$};
\draw (EE) -- (E1) (EE) -- (E2) (E1) -- (K) (E2) -- (K) (K) -- (F);
\node[font=\scriptsize,cold] at (0,-1.6) {все расширения~--- Галуа над $F$};
\end{scope}
\begin{scope}[xshift=7cm]
\node[cnew] (G) at (0,2.4) {$\Gal(E_1E_2/F)$};
\node[cold] (G1) at (-1.7,1.2) {$\Gal(E_1/F)$};
\node[cold] (G2) at (1.7,1.2) {$\Gal(E_2/F)$};
\node[cnew] (G0) at (0,0) {$\Gal(E_1\cap E_2/F)$};
\draw[->>] (G) -- node[above left,font=\scriptsize]{$\sigma\mapsto\sigma|_{E_1}$} (G1);
\draw[->>] (G) -- node[above right,font=\scriptsize]{$\sigma\mapsto\sigma|_{E_2}$} (G2);
\draw[->>] (G1) -- node[below left,font=\scriptsize]{$\pi_1$} (G0);
\draw[->>] (G2) -- node[below right,font=\scriptsize]{$\pi_2$} (G0);
\node[font=\scriptsize,cold] at (0,-1.0) {расслоённое произведение (декартов квадрат)};
\end{scope}
\end{tikzpicture}
\end{center}
\figcap{Предложение \rf{3.21}: ромбу полей отвечает ромб групп Галуа со стрелками ограничения; группа Галуа
композита~--- расслоённое произведение групп Галуа полей $E_1$ и $E_2$ над группой Галуа их пересечения.}
\end{keepfig}

\begin{proof}
\emph{$E_1\cap E_2$~--- Галуа.} Пусть $a\in E_1\cap E_2$ и $f$~--- его минимальный многочлен над $F$. Так как $E_1/F$ нормально и
сепарабельно, $f$ имеет $\deg f$ различных корней в $E_1$; так же~--- в $E_2$. В поле $E_1E_2$ у $f$ не более $\deg f$
корней, поэтому оба набора корней совпадают и лежат в $E_1\cap E_2$. Значит, каждый неприводимый многочлен из $F[X]$,
имеющий корень в $E_1\cap E_2$, имеет там $\deg f$ различных корней, т.\,е.\ $E_1\cap E_2$ нормально и сепарабельно над $F$;
оно конечно (лежит в $E_1$), и по определению это расширение Галуа.

\emph{$E_1E_2$~--- Галуа.} Пусть $E_i$~--- поле разложения сепарабельного $f_i\in F[X]$ (\rf{3.10}), и $f$~--- произведение
различных унитарных неприводимых множителей $f_1$ и $f_2$ (наименьшее общее кратное $f_1$ и $f_2$, если они унитарны); он сепарабелен (различные неприводимые многочлены не имеют
общих корней, \rf{2.17}). Поле $E_1E_2$ порождено над $F$ корнями $f_1$ и $f_2$, т.\,е.\ корнями $f$, и $f$ раскладывается
в нём на линейные множители; значит, $E_1E_2$~--- поле разложения сепарабельного многочлена $f$ и расширение Галуа.

\emph{Гомоморфизм.} Каждый $\sigma\in\Gal(E_1E_2/F)$ переставляет корни $f_i$, поэтому $\sigma E_i=E_i$, и
$\varphi(\sigma)=(\sigma|_{E_1},\sigma|_{E_2})$~--- гомоморфизм в $\Gal(E_1/F)\times\Gal(E_2/F)$. Он инъективен: если $\sigma$
тождествен на $E_1$ и на $E_2$, то и на $E_1E_2$. Его образ лежит в $H$: оба ограничения совпадают с $\sigma|_{E_1\cap E_2}$.

\emph{Образ равен $H$.} Подсчитаем порядок $H$. По \rf{3.17}(d) для расширения Галуа $E_2/F$ и его нормального подрасширения
$K=E_1\cap E_2$ ограничение $\pi_2\colon\Gal(E_2/F)\to\Gal(K/F)$~--- сюръективный гомоморфизм с ядром $\Gal(E_2/K)$; поэтому
каждый элемент $\Gal(K/F)$ имеет ровно $|\Gal(E_2/K)|=[E_2:K]$ прообразов. Элемент $(\sigma_1,\sigma_2)\in H$ однозначно
задаётся выбором $\sigma_1\in\Gal(E_1/F)$ ($[E_1:F]$ возможностей) и затем выбором $\sigma_2$ среди прообразов
$\sigma_1|_K$ ($[E_2:K]$ возможностей); разные выборы дают разные пары. Значит,
\[
|H|=[E_1:F]\,[E_2:K]=\frac{[E_1:F]\,[E_2:F]}{[K:F]}=[E_1E_2:F]
\]
(второе равенство~--- \rf{1.20}, третье~--- \rf{3.20}). Итак, инъективный гомоморфизм $\varphi$ отображает группу порядка
$[E_1E_2:F]$ в группу $H$ того же порядка и потому является изоморфизмом на $H$.
\end{proof}

\begin{corollary}[3.21a]
В условиях \rf{3.21} $\Gal(E_1E_2/F)\cong\Gal(E_1/F)\times\Gal(E_2/F)$ (посредством ограничений) тогда и только тогда, когда
$E_1\cap E_2=F$; в этом случае $[E_1E_2:F]=[E_1:F]\,[E_2:F]$.
\end{corollary}

\begin{proof}
Если $E_1\cap E_2=F$, то условие на пары пусто и $H$~--- всё произведение. Обратно, если $H$~--- всё произведение, то
$|H|=[E_1:F][E_2:F]$, а по доказанному $|H|=[E_1:F][E_2:F]/[E_1\cap E_2:F]$, откуда $E_1\cap E_2=F$.
\end{proof}

Например, $\Q[\sqrt2]\cap\Q[\sqrt3]=\Q$ (упражнение 2-1), поэтому $\Gal(\Q[\sqrt2,\sqrt3]/\Q)\cong\Z/2\times\Z/2$; так же
$\Q[\zeta_m]\cap\Q[\zeta_n]=\Q$ при $\gcd(m,n)=1$, и $\Gal(\Q[\zeta_{mn}]/\Q)\cong\Gal(\Q[\zeta_m]/\Q)\times\Gal(\Q[\zeta_n]/\Q)$
(\rf{5.11c}). Если же $E_1=\Q[\sqrt[3]2,\omega]$ (группа $S_3$) и $E_2=\Q[\sqrt{-3}]=\Q[\omega]\subset E_1$, то
$E_1\cap E_2=E_2$, $E_1E_2=E_1$, и $H=\{(\sigma,\sigma|_{E_2})\}\cong S_3$~--- <<график>> гомоморфизма знака
$S_3\to\Z/2$, а не всё произведение $S_3\times\Z/2$. Для группы Галуа приводимого многочлена $f=f_1f_2$ это означает, что
$G_f$~--- расслоённое произведение $G_{f_1}$ и $G_{f_2}$ над группой пересечения их полей разложения (подробнее~---
\rf{3.27e}).

\subsection*{Композит как склейка: связь с кодекартовыми квадратами}

В \rf{2.15a}--\rf{2.15b} мы видели, что у двух расширений $E$ и $L$ поля $K$, вообще говоря, нет <<универсальной склейки>>
(кодекартова квадрата): композитов $E$ и $L$ может быть несколько, и они не сводятся один к другому. Оказывается, если $E$~---
расширение Галуа, а склеивание ведётся над пересечением $K=E\cap L$, то композит $EL$ является именно такой
универсальной склейкой; группы Галуа превращают этот кодекартов квадрат полей в декартов квадрат групп~--- расслоённое
произведение \rf{3.21}.

\begin{proposition}[3.21b: склейка над пересечением]
Пусть $E$ и $L$~--- расширения $F$ в общем поле, $E/F$~--- расширение Галуа, $K=E\cap L$. Для любого поля $M$ и любых
гомоморфизмов $\tau_1\colon E\to M$, $\tau_2\colon L\to M$, совпадающих на $K$, существует единственный гомоморфизм
$\tau\colon EL\to M$ с $\tau|_E=\tau_1$ и $\tau|_L=\tau_2$. Иначе говоря:
\begin{enumerate}[label=\textup{(\alph*)}]
\item квадрат из полей $K\subset E$, $K\subset L$, $E\subset EL\supset L$ кодекартов в смысле \rf{2.15b} (с $K$ в роли
основного поля);
\item для любого расширения $M$ поля $F$ ограничения дают биекцию
\[
\Hom_F(EL,M)\ \xrightarrow{\ \sim\ }\ \Hom_F(E,M)\times_{\Hom_F(K,M)}\Hom_F(L,M),\qquad \tau\mapsto(\tau|_E,\tau|_L).
\]
\end{enumerate}
\end{proposition}

\begin{proof}
\emph{Единственность.} Поле $EL$ порождено $E\cup L$: каждый его элемент~--- частное двух сумм произведений элементов $E$
и $L$, и значение $\tau$ на нём определяется значениями на $E$ и на $L$.

\emph{Существование.} Пусть $E$~--- поле разложения над $F$ сепарабельного многочлена $f\in F[X]$, $\alpha_1,\dots,\alpha_n$~---
его корни в $E$. Гомоморфизмы $\tau_1,\tau_2$ совпадают на $F\subset K$; обозначим $\tau_0=\tau_1|_F=\tau_2|_F$. Применяя $\tau_1$ к разложению $f=\prod(X-\alpha_i)$ (считаем $f$ унитарным), получаем, что
$\tau_0f=\prod(X-\tau_1\alpha_i)$ раскладывается в $M$ на различные линейные множители. Поле $EL=L[\alpha_1,\dots,\alpha_n]$
порождено над $L$ корнями $f$, и $\tau_2f=\tau_0f$ раскладывается в $M$; по \rf{2.12a} гомоморфизм $\tau_2$ продолжается до
гомоморфизма $\tau'\colon EL\to M$. Теперь $\tau'|_E$ и $\tau_1$~--- два гомоморфизма $E\to M$, совпадающие на $K$
(оба равны там $\tau_2$). Каждый из них переводит множество $\{\alpha_i\}$ корней $f$ биективно на множество
$\{\tau_1\alpha_i\}$ корней $\tau_0f$ в $M$, поэтому их образы совпадают: $\tau'(E)=\tau_1(E)=\tau_0(F)[\tau_1\alpha_1,\dots,\tau_1\alpha_n]$.
Значит, $\rho=(\tau'|_E)^{-1}\circ\tau_1$~--- автоморфизм поля $E$, тождественный на $K$, т.\,е.\ $\rho\in\Gal(E/K)$. По
\rf{3.19} найдётся $\tilde\rho\in\Gal(EL/L)$ с $\tilde\rho|_E=\rho$. Положим $\tau=\tau'\circ\tilde\rho$. Тогда
$\tau|_L=\tau'|_L=\tau_2$ (так как $\tilde\rho$ тождествен на $L$), а $\tau|_E=\tau'|_E\circ\rho=\tau_1$.

Утверждения (a) и (b)~--- переформулировки: в (a) поле $M$ рассматривается как расширение $K$ посредством общего
ограничения $\tau_1|_K=\tau_2|_K$, а в (b) пары $F$-гомоморфизмов, совпадающих на $K$,~--- это элементы расслоённого
произведения множеств.
\end{proof}

\begin{keepfig}
\begin{center}
\begin{tikzcd}[column sep=small,row sep=normal]
 & & M & & & & \\
 & EL\arrow[ur,cnew,dashed,"\exists!\,\tau" description] & & & & \Gal(EL/F)\arrow[dl]\arrow[dr] & \\
E\arrow[ur,hook]\arrow[uurr,cold,bend left=30,"\tau_1"] & & L\arrow[ul,hook]\arrow[uu,cold,"\tau_2"'] & \qquad\leadsto\qquad & \Gal(E/F)\arrow[dr] & & \Gal(L/F)\arrow[dl] \\
 & K=E\cap L\arrow[ul,hook]\arrow[ur,hook] & & & & \Gal(K/F) &
\end{tikzcd}
\end{center}
\figcap{Слева: композит $EL$~--- кодекартов квадрат (универсальная склейка $E$ и $L$ над $K=E\cap L$), если $E/F$~---
расширение Галуа (\rf{3.21b}). Справа: при $L/F$ тоже Галуа соответствие Галуа обращает стрелки, и тот же квадрат
превращается в декартов квадрат групп~--- расслоённое произведение (\rf{3.21}).}
\end{keepfig}

\begin{remark}[3.21c]
(1) \emph{Второе доказательство \rf{3.21}.} Пусть $E_1,E_2$~--- расширения Галуа поля $F$ и $M=E_1E_2$. Каждый
$F$-гомоморфизм $E_i\to E_1E_2$ переводит корни многочлена $f_i$ (полем разложения которого является $E_i$) в корни $f_i$,
т.\,е.\ в $E_i$, и является автоморфизмом $E_i$ (инъективное линейное отображение конечномерного пространства в себя);
поэтому $\Hom_F(E_i,E_1E_2)=\Gal(E_i/F)$, и так же $\Hom_F(E_1E_2,E_1E_2)=\Gal(E_1E_2/F)$,
$\Hom_F(E_1\cap E_2,E_1E_2)=\Gal(E_1\cap E_2/F)$. Биекция \rf{3.21b}(b) превращается в изоморфизм \rf{3.21}~--- без подсчёта
порядков. Именно так \rf{3.21} переносится на бесконечные расширения Галуа (\rf{7.15}--\rf{7.16b}), где считать порядки нельзя.

(2) \emph{Связь с критерием \rf{2.15b}.} Если $E=K[\gamma]$ и $g$~--- минимальный многочлен $\gamma$ над $K$, то $EL=L[\gamma]$, и
равенство $[EL:L]=[E:K]$ из \rf{3.19} означает, что $g$ остаётся неприводимым над $L$, т.\,е.\ выполнено условие $r=1$ из
\rf{2.15b}; это согласуется с (a).

(3) \emph{Без условия Галуа склейки нет.} В примере \rf{3.19a}(3) ($E=\Q[\sqrt[3]2]$, $L=\Q[\omega\sqrt[3]2]$, $K=\Q$) над $L$
имеем $X^3-2=(X-\omega\sqrt[3]2)(X^2+\omega\sqrt[3]2\,X+\omega^2\sqrt[3]4)$~--- два неприводимых множителя ($\sqrt[3]2\notin L$), и по \rf{2.15b}
кодекартова квадрата нет. Конкретно: $\Q$-гомоморфизмы $\tau_1\colon\sqrt[3]2\mapsto\sqrt[3]2$ и $\tau_2\colon\omega\sqrt[3]2\mapsto\sqrt[3]2$
(в $M=\R$) совпадают на $K=\Q$, но не склеиваются: склейка перевела бы $\omega=\omega\sqrt[3]2/\sqrt[3]2$ в $1$. Поля $E$ и $L$
можно <<склеить>> по-разному: в композит $\Q[\sqrt[3]2,\omega]$ степени $6$ или, отождествив $L$ с $E$, в поле степени $3$.

(4) \emph{Склейки поля Галуа с самим собой.} Пусть $E/K$~--- конечное расширение Галуа с группой $G\ne1$ и $E=K[\gamma]$.
Минимальный многочлен $g$ элемента $\gamma$ раскладывается над $E$ на $|G|\ge2$ различных линейных множителей
$X-\sigma\gamma$, поэтому по \rf{2.15b} кодекартова квадрата для пары $K\subset E$, $K\subset E$ нет. Но все <<склейки>> легко
перечислить. Пусть $\tau_1,\tau_2\colon E\to M$ совпадают на $K$. Поле $\tau_i(E)$ порождено над $\tau_1(K)$ образами
$\deg g$ корней $g$ в $E$, т.\,е.\ всеми корнями многочлена $\tau_1g=\tau_2g$ в $M$ (больше $\deg g$ корней у него нет);
поэтому $\tau_1(E)=\tau_2(E)$, и $\sigma=\tau_1^{-1}\circ\tau_2$~--- автоморфизм $E$ над $K$, т.\,е.\ $\sigma\in G$. Значит,
$(\tau_1,\tau_2)=(\tau_1\circ\id,\tau_1\circ\sigma)$: каждая склейка получается из пары $(\id,\sigma)\colon E\rightrightarrows E$
вложением $\tau_1$, и $\sigma$ определён однозначно. Итак, склейки $E$ с собой над $K$~--- это само поле $E$, склеенное
посредством одного из $|G|$ автоморфизмов. На языке \rf{2.15a} это разложение
$E\otimes_KE\cong E[X]/(g)\cong\prod_{\sigma\in G}E[X]/(X-\sigma\gamma)\cong\prod_{\sigma\in G}E$.
\end{remark}
